% ===== §8 The Impossibility of Writing as Self-Verification =====
\section{The Impossibility of Writing as the Self-Verification of the Thesis}
\label{p19:sec:unwritability}

\noindent
\emph{The hard core (against ``this is mere sentiment'').} A chapter saying ``I cannot write it'' is, without argument, a sentimental self-congratulation. This chapter's load is to show that the unwritability of this paper's object is \emph{not a failure of the writer but structural}---a direct consequence of the symbolic--real gap of \xref{p19:sub:gen-drive}. Writing is an operation of the symbolic; the still afternoon holds a real remainder no symbol can exhaust; therefore ``it cannot be written to the end'' is necessary, not rhetorical. From this, the unwritability \emph{verifies} the paper's core claim (the gap unclosable, happiness holding a remainder no symbol exhausts). This is argument, not sentiment.

\subsection{Returning to the quarry of the opening}
\label{p19:sub:unwrit-quarry}

Return to the quarry laid down at the opening (the delineation): why the harder one tries to write that afternoon clearly, the more it resists. This chapter raises that ``resistance'' from a fact of writing-experience to a phenomenon \emph{predicted} by the foregoing structure.

\subsection{Unwritability is structural, not a matter of ability}
\label{p19:sub:unwrit-structural}

Writing is an operation of the symbolic: to capture, fix, and transmit by the order of signifiers. But (per \xref{p19:sub:gen-drive}) the symbolic's capture of the real always leaves a remainder; that afternoon, as a lived situation under conferral of meaning, holds a real remainder no sequence of signifiers can exhaust (the remainder marked by \emph{objet a}). Hence: \emph{however precise the writing, there is always a remainder that does not enter---this is not the author's limit but the limit of the symbolic--real gap, and it is structural.} ``Cannot write that afternoon to the end'' and ``the symbol cannot fill the real'' are one proposition. Following \xref{p19:sub:gen-manifold}: writing (the symbolic) too can only unfold on the background manifold, and is never the written life itself; between text and life lie the manifold and the difference of meaning. The paper is a map, not the terrain.

\subsection{This impossibility verifies the thesis (the self-referential closure)}
\label{p19:sub:unwrit-verify}

One of the paper's core claims: happiness, as the fullness of a generative relation, holds a remainder no symbol exhausts (the gap abiding in the mode of fullness, \xref{p19:sub:gen-twomodes}); it is therefore a genuine emergence, hermeneutic, not to be settled by any form or statement. Then: \emph{could this paper write that afternoon exhaustively, it would thereby falsify its own core claim} (proving happiness exhaustible by the symbol, settleable). Just because it cannot be exhausted, the core claim is borne out. \emph{Unwritability is the performative proof of the claim's truth.}

The convergence with the series' self-referential gesture:
\begin{itemize}
  \item Paper~IX \parencite{paperix}: the poem refuses to settle its own sublimation---its form \emph{is} its claim.
  \item Paper~XIII \parencite{paperxiii}: the dedication enacts the very theory it states---the performative \emph{is} the argument.
  \item This chapter: the paper cannot write its object to the end, which is just what proves the object inexhaustible by writing---the failure \emph{is} the evidence of success.
  \item (And \xref{p19:sec:dynamics}, ``insufficiency-cum-bearing,'' is the formal-dress version of the same gesture: exhausting the form to display the form's not-reaching.)
\end{itemize}
The three (four) are one self-referential gesture in successive instances.

\subsection{But this chapter does not slide into the unsayable or the mystical (self-brake one)}
\label{p19:sub:unwrit-brake1}

A guard: against sliding from ``cannot be written to the end'' to ``therefore unsayable, therefore mystical, therefore nothing need be argued.'' No. ``Cannot be written to the end'' $\neq$ ``unsayable'': this paper has said a great deal (eight chapters)---it only does not say \emph{exhaustively}. The existence of a remainder does not annul the validity of what has been said (just as the paths one cannot walk on the manifold do not annul the paths one can, \xref{p19:sub:gen-manifold}). The saying remains under constraint (the four criteria of \xref{p19:sub:gen-methodology}, the double constraint of \xref{p19:sub:gen-manifold}): unwritability is an honest marking of the boundary of saying, not a use of ``the unsayable'' to abdicate the responsibility to argue. This chapter is the surveying of a boundary, not a forfeiture of argument.

\subsection{Nor into self-congratulation (self-brake two)}
\label{p19:sub:unwrit-brake2}

The chapter's greatest risk is to turn ``I cannot write it to the end'' into ``I am too profound to write it to the end'' (the same slope as the ``unfalsifiability as virtue'' guarded against in \xref{p19:sub:gen-methodology}). The brake: unwritability is a structural property \emph{of the object} (the symbolic--real gap), not a special profundity \emph{of the author}. \emph{Anyone} writing \emph{any} lived situation cannot write it to the end. The chapter asserts a general structure of symbol and real that happens to fall on this paper's object---\emph{not} the author's self-coronation.

\subsection{Close: the fitting end is to lay down the pen}
\label{p19:sub:unwrit-close}

Since the third silence (\xref{p19:sub:silence-three}) is ``no need to say,'' and this paper has said what can be said to its boundary, the fitting end of this paper is not one more conclusion (a conclusion is a settlement, and settlement is the death of the good cycle, \xref{p19:sec:generativity}) but the \emph{laying-down of the pen}. To lay down the pen is silence; silence is fulfilment---the paper here instantiates what it argues: that the fulfilment of the plain happiness need not be stated, and so, at the boundary, it stops, adding no further signifier. An opening is left for the envoi.
